\documentclass[11pt,letterpaper]{article}
\usepackage[margin=0.83in,headheight=15pt]{geometry}
\usepackage{fontspec}
\setmainfont{Bitstream Charter}
\setsansfont{DejaVu Sans}
\setmonofont{DejaVu Sans Mono}[Scale=0.82]
\usepackage{amsmath,amssymb,mathtools,booktabs,longtable,array,enumitem}
\usepackage{xcolor,fancyhdr,titlesec,fvextra}
\usepackage{hyperref,xurl}
\definecolor{ink}{HTML}{17324D}
\definecolor{accent}{HTML}{167568}
\definecolor{muted}{HTML}{546477}
\hypersetup{colorlinks=true,linkcolor=ink,urlcolor=accent,pdftitle={Four Further Lean Proof Simplifications},pdfauthor={Prepared for Vladimir Reshetnikov}}
\titleformat{\section}{\Large\sffamily\bfseries\color{ink}}{\thesection}{0.65em}{}
\titleformat{\subsection}{\large\sffamily\bfseries\color{ink}}{\thesubsection}{0.65em}{}
\pagestyle{fancy}\fancyhf{}
\fancyhead[L]{\small\sffamily Further Lean proof simplification}
\fancyhead[R]{\small\sffamily 26 October 2026}
\fancyfoot[C]{\small\thepage}
\setlength{\parindent}{0pt}\setlength{\parskip}{6pt}
\setlist{nosep,leftmargin=1.6em}
\newcommand{\code}[1]{\nolinkurl{#1}}
\urlstyle{tt}
\setlength{\emergencystretch}{2em}
\newcommand{\repo}{https://github.com/VladimirReshetnikov/ProveIt}
\newcommand{\pin}{7df43cdfacb490cd3ca3acd23e39027efebf8801}
\newcommand{\src}[2]{\href{\repo/blob/\pin/Computability/DiophantineConsistency/Lean/DiophantineConsistency/#1.lean}{\code{#1}: #2}}
\newcolumntype{L}[1]{>{\raggedright\arraybackslash}p{#1}}
\newcommand{\N}{\mathbb N}\newcommand{\Z}{\mathbb Z}
\newcommand{\SZ}{\operatorname{SZ}}
\DefineVerbatimEnvironment{Lean}{Verbatim}{fontsize=\small,samepage=true,breaklines=true,breakanywhere=true,frame=single,rulecolor=\color{muted},framesep=6pt}
\begin{document}
\begin{titlepage}
\vspace*{0.5in}
{\Large\sffamily\color{accent} PROVEIT / DIOPHANTINE CONSISTENCY\par}
\vspace{0.28in}
{\Huge\sffamily\bfseries\color{ink} Four Further Lean\\[4pt] Proof Simplifications\par}
\vspace{0.25in}
{\Large Polynomial identities, list sums,\\assignment restriction, and exponent divisibility\par}
\vspace{0.35in}
{\large Prepared for Vladimir Reshetnikov\par}
26 October 2026
\vspace{0.35in}
\hrule
\vspace{0.18in}
\textbf{Main result.} A checked transport theorem lets equality of integer polynomials
justify equality-of-splits semantics over every additively cancellative commutative
semiring. It supports arbitrary internal naturals without identifying them with
external natural numbers.

\textbf{Concrete companions.} A complete sparse-coefficient module refactor removes
four of five list inductions and ten source lines. A smaller exponent proof removes
four lines. An environment proof proposes twelve fewer lines; its abstract
mechanism is checked, while its full object-model integration remains uncompiled.

\textbf{Exclusion boundary.} Earlier proposals are treated as already applied,
including optional and uncompiled ones. The audit covers 23 placed packets and
five incoming ZIPs. A late packet caused a coprime-divisibility candidate to be
withdrawn.
\vfill
{\small\color{muted}
Source and compiler baseline: \code{89226e3921bc}.\\
Closing novelty review: \code{7df43cdfacb4}.\\
Lean 4.32.0, with the repository's pinned Mathlib.\\
Source, patches, compiler receipts, and exclusion evidence are bundled.\\
No whole-project build or global source-size reduction is claimed.}
\end{titlepage}
\begingroup
\small
\setlength{\parskip}{0pt}
\tableofcontents
\endgroup
\newpage
\section{Findings and recommended order}

The remaining opportunities are narrower than the earlier broad refactors.
After treating the prior scanner, compiler, coefficient, geometry, and
object-language interfaces as available, four distinct mechanisms survive
the exclusion audit. The first is a reusable semantic theorem; the other
three simplify specific maintained proofs.

\begin{longtable}{L{0.06\linewidth}L{0.43\linewidth}L{0.40\linewidth}}
\toprule ID & Change and target & Evidence and cost\\ \midrule
P1 & Transfer integer polynomial identities to natural split semantics;
target algebraic row proofs in the object layer. &
Core and literal selection-row client checked. New shared lemmas with an
up-front proof cost; no measured project-wide saving.\\
P2 & View the existing sparse-list interpreter as a monomial sum;
target \code{SyntaxCoefficients}. &
Complete original and candidate modules checked. 120 to 110 physical lines;
97 to 87 nonblank lines.\\
P3 & Replace manual environment extension by restriction surjectivity;
target \code{sat_noZeroForm_iff}. &
Generic mechanism checked. Production declaration proposed at 20 to 8 lines;
full object-model integration uncompiled.\\
P4 & Prove the residual exponent is zero from an order obstruction;
target \code{width_dvd_exponent}. &
Exact old and new declarations checked, with identical signatures and axiom
sets. 12 to 8 lines. Optional prime-radix API checked separately.\\
\bottomrule
\end{longtable}

For immediate adoption, begin with P2: its entire touched module was checked
against the freshly rebuilt current representation module. P4 is also small,
but its enclosing module and clients still need an integration build.
P3 is a reviewable proof replacement whose final object-language elaboration
is explicitly outstanding. P1 deserves a pilot across several algebraic rows
before deciding how much shared infrastructure to retain.

The local changes sum to 26 \emph{proposed} removed lines across P2--P4.
That total mixes validation scopes: ten lines come from a complete checked
module, four from checked exact declarations, and twelve from a proposed
production integration. P1 adds new lemmas, whose cost must not be hidden
inside that total. There is no claim that the packet makes the repository
smaller immediately.

\subsection{Source baseline and closing refresh}

The initial GitHub main branch was
\href{\repo/commit/89226e3921bce3e8949fadd12d3efb6914c7630c}{\code{89226e3921bce3e8949fadd12d3efb6914c7630c}}.
All 653 blobs in its maintained Lean subtree were matched to the GitHub tree
using Git blob SHA-1. The detailed inspection and candidate construction use
that snapshot.

A bounded final comparison through
\href{\repo/commit/7df43cdfacb490cd3ca3acd23e39027efebf8801}{\code{7df43cdfacb490cd3ca3acd23e39027efebf8801}}
found eleven added Lean modules, aggregate imports, research reviews, and
one more incoming ZIP. None of the existing target modules changed.
The added report was read and included in the exclusion audit.
The comparison receipt is bundled.

This is a snapshot review, not a claim about every later state of a moving
branch. The patches identify their exact original bytes.

\section{What has already been proposed}

The exclusion baseline is the union of all proposals in
\href{\repo/tree/\pin/Computability/DiophantineConsistency/research/lean-proof-simplification}
{the 23 placed packets a--w} and the five ZIPs in
\href{\repo/tree/\pin/docs/incoming}{the refreshed incoming directory}.
The audit reads main reports and substantive supplementary notes, including
optional designs, rather than only adoption records. The incoming README
was used for routing and relevant searches; unrelated historical manuscripts
transitively linked from it were not all reread.

The following families receive no new credit here:

\begin{itemize}
\item Optional-list lookup, exact decoder graphs, scanner iteration,
validated histories, Boolean reflection, and rule-schedule interfaces.
\item Assignment naturality, witness-block layouts, finite residual sums,
generic degree and resource bounds, circuit generation, and topological
evaluation invariants.
\item Coefficient envelopes, finite sparse reflection, coefficient fibres,
row operations, and the earlier semiring gate interpreters.
\item Canonical carries, cyclic packing, aligned-cell subgroups, helical
periods and coordinates, fixed Pell anchors, and coherent word windows.
\item Internal definable induction, unique-existence and graph fusion,
shared model bridges, internal order constructors, and sentence transport.
\end{itemize}

\subsection{Nearest-precedent comparison}

\begin{longtable}{L{0.06\linewidth}L{0.40\linewidth}L{0.43\linewidth}}
\toprule ID & Closest excluded mechanism & Distinct remaining obligation\\ \midrule
P1 & Standard-natural specialization of semiring gates; reuse of internal
evaluators. &
Reflect integer-polynomial equality to a natural cross-sum identity and
preserve signed zero semantics in arbitrary cancellative commutative semirings.\\
P2 & Finite row coefficient tables and sparse row reflection. &
Interpret the existing executable duplicate-preserving list expansion via
one sum theorem, replacing several list inductions.\\
P3 & Arity-polymorphic closed-sentence transport. &
Simplify the environment-extension proof inside the semantic equivalence
consumed by those wrappers.\\
P4 & Coherent source/window APIs assuming recovered powers. &
Simplify the exact exponent-remainder proof; optionally recover a block
power directly from arithmetic divisibility certificates.\\
\bottomrule
\end{longtable}

Several apparent opportunities were rejected. Root-constructor minima for
\code{AnnotatedClosure.formCode_ge} are already in the proof-refinements ZIP.
Periodic-orbit and lattice descriptions of helices occur in several packets.
Signed-prefix borrowing occurs in the last additional interfaces report.
That last report also contains the exact
\code{RatioReparameterization.combined_dvd_iff} simplification first considered
during this session; its candidate code and claimed savings are omitted.

The final ledger contains 115 inherited contract rows and 61 later overlapping
commitments, with 68 report or section-file records across 28 packets.
These are audit records, not 176 independent mathematical ideas.
The appropriate novelty claim is that no matching mechanism was found in
this explicitly audited corpus, not global originality.

\section{P1: transport polynomial identities into internal arithmetic}

\subsection{The proof obligation left by a semiring evaluator}

The maintained development has three relevant layers.
\src{Representation}{\code{Expr}} stores integer arithmetic syntax and its
ordinary integer polynomial.
\src{NaturalExpression}{\code{splitNatural}} translates a signed expression
into two natural-coefficient expressions.
\src{GatesSemiring}{\code{SemiringZero}} interprets vanishing as equality of
their natural values.

If $e$ has split $(e^+,e^-)$, write
\[
 \SZ_R(e,v)\quad:\Longleftrightarrow\quad
 [\![ e^+]\!]_R(v)=[\![ e^-]\!]_R(v).
\]
Subtraction need not exist on $R$. An existing semiring evaluator does not
automatically permit replacing $e$ by another expression $f$ with the same
integer polynomial: their positive and negative syntax pieces can look quite
different. Repeatedly unfolding \code{splitNatural} exposes this implementation
to every client.

The proposed interface is
\[
 \boxed{\quad e.\mathrm{toPolynomial}=f.\mathrm{toPolynomial}
 \ \Longrightarrow\
 \bigl(\SZ_R(e,v)\leftrightarrow\SZ_R(f,v)\bigr).\quad}
\]
It holds for every additively cancellative commutative semiring $R$.
A proof can establish a finite algebraic normalization in an ordinary integer
polynomial ring and then reuse it on any appropriate internal-natural carrier.
The stored expression itself remains unchanged.

\subsection{Proof by a natural cross-sum identity}

Let $P_e,N_e,P_f,N_f\in\N[\mathbf X]$ be the natural-coefficient polynomials
of the four split expressions. Let
\[
 \iota:\N[\mathbf X]\longrightarrow\Z[\mathbf X]
\]
be the homomorphism induced by the natural-to-integer coefficient map.

One structural induction proves
\[
 \iota(P_e)-\iota(N_e)=e.\mathrm{toPolynomial}.
 \tag{1}
\]
The constant case is the integer identity
$a.\mathrm{toNat}-(-a).\mathrm{toNat}=a$.
For multiplication, the positive split is $P_eP_f+N_eN_f$ and the negative
split is $P_eN_f+N_eP_f$, so the invariant is the product-of-differences
identity. No semantic model occurs in this induction.

If $e$ and $f$ have equal integer polynomials, (1) gives
\[
 \iota(P_e)-\iota(N_e)=\iota(P_f)-\iota(N_f),
\]
hence
\[
 \iota(P_e+N_f)=\iota(P_f+N_e).
\]
Injectivity of the coefficient map, and therefore its multivariate-polynomial
map, reflects this to
\[
             P_e+N_f=P_f+N_e\quad\text{in }\N[\mathbf X].
 \tag{2}
\]
Evaluate (2) through the homomorphism induced by \code{Nat.castRingHom}~\code{R}:
\[
             p_e+n_f=p_f+n_e.
 \tag{3}
\]
If $p_e=n_e$, substitute and cancel a common additive term to obtain
$p_f=n_f$. The converse is symmetric.

The proof does not embed the semantic carrier in a ring or construct its
Grothendieck completion. It translates a finite coefficient identity once
and retains the original semiring semantics.

\subsection{Checked Lean interface}

The archive's \code{PolynomialZeroTransport.lean} proves four facts about
the actual repository syntax:

\begin{Lean}
map_splitNatural_sub
splitNatural_cross_sum
splitNatural_eval_cross_sum
splitNatural_zero_iff_of_polynomial_eq
\end{Lean}

Its main theorem is:

\begin{Lean}
theorem splitNatural_zero_iff_of_polynomial_eq
    [CommSemiring R] [IsCancelAdd R]
    (e f : Expr σ)
    (h : e.toPolynomial = f.toPolynomial) (v : σ → R) :
    (e.splitNatural.1.evalSemiring v =
       e.splitNatural.2.evalSemiring v) ↔
    (f.splitNatural.1.evalSemiring v =
       f.splitNatural.2.evalSemiring v)
\end{Lean}

A production wrapper using the existing \code{SemiringZero} name is
definitionally the same statement. The core belongs below the gate and
object-model layers to avoid an import cycle. Its checked file imports
\code{NaturalExpressionSemiring}, not the deep object-model closure.

\subsection{Concrete client: the conditional selection row}

The literal row in
\src{ObjectConditional}{\code{literal_selection}} is
\[
       e=o-\bigl(bz+(1-b)p\bigr).
\]
The object-level goal is the balanced equality
\[
       o+bp=bz+p.
 \tag{4}
\]
Normalize to
\[
       f=(o+bp)-(bz+p).
\]
The equality of their ordinary integer polynomials follows by \code{ring}.
P1 transfers the zero predicate from $e$ to $f$; the latter has immediate
split semantics (4).

The checked client leaves all four variable indices abstract and then
instantiates the exact slots used by \code{ObjectConditional}: output $0$,
bit $4$, zero-branch $2$, and positive-branch $3$, with an empty register
block. This checks the actual row shape, not a numerical analogue.

This row already has a short direct proof. It demonstrates the interface,
not a net saving by itself. The useful adoption experiment is a group of
algebraically normalized rows in the zero-test, conditional, pairing, and
division syntax layers. Purely structural gates can retain their current
cheap specialized lemmas.

\subsection{The hypotheses are substantive}

Additive cancellation cannot be removed. In the Boolean semiring, with
addition OR and multiplication AND, consider
\[
              e=1,\qquad f=2-1.
\]
Their integer polynomials are both $1$. But $\SZ(e)$ asks $1=0$, which is
false, while $\SZ(f)$ asks $2=1$, which is true in this semiring.
The supplied client file constructs the Boolean semiring and checks both
evaluations. This is an explicit negative control on the theorem's scope.

Commutative multiplication also matters. Equality in a commutative
polynomial ring cannot justify $xy=yx$ in a noncommutative matrix semiring.

Conversely, P1 requires no order, characteristic-zero hypothesis, absence of
zero divisors, or multiplicative cancellation. The intended internal-natural
application uses the repository's existing semiring and additive-cancellation
structure. It never identifies a possibly nonstandard model with external
\code{Nat}.

\subsection{Adoption cost and limits}

The supplied core has 72 physical lines, including comments and three axiom
print commands. It is an added interface, whose whole cost belongs in the
accounting. Adopt it when several clients can state their desired balanced
equation and discharge only a polynomial identity, with less dependence on
the recursive split implementation.

Retain the literal expression, polynomial, witness counts, degree
certificates, and all object-model assumptions. The checked core and client
prove the semantic mechanism; they do not certify a full rebuild of the
internal ZFC proof chain. No project-wide line saving is claimed for P1.

\section{P2: give sparse-list expansion one monomial-sum view}

\subsection{The repeated reasoning}

\src{SyntaxCoefficients}{the interpretation lemmas} keeps sparse expansions
as executable lists. A term consists of an exponent vector and an integer
coefficient. Multiplication distributes through \code{flatMap} and
\code{map}, and duplicate monomials remain in the list.

The recursive \code{interpret} function denotes the corresponding ordinary
polynomial sum. The original proof layer independently inducts over lists
for append, coefficient negation, multiplication by one term, full
multiplication, and coefficient extraction. These are all algebraic facts
about list sums.

Keep the representation and every executable definition. Add just one
private view:
\[
 \operatorname{interpret}(L)=
 \sum_{t\in L}\operatorname{monomial}(\operatorname{exponent}(t_1),t_2),
 \tag{5}
\]
where the sum is an occurrence-sensitive list sum, not a finite set sum.

\subsection{The new proof structure}

Prove (5) once by list induction. Then use the following interfaces:

\begin{itemize}
\item Append: \code{List.map_append} and \code{List.sum_append}.
\item Negation: the additive-map law for monomial coefficients and
\code{List.sum_neg}.
\item Multiplication: flatten a list of lists, then use
\code{List.map_flatten}, \code{List.sum_flatten}, and distribution over sums.
\item Coefficient extraction: apply the bundled
\code{MvPolynomial.coeffAddMonoidHom} to the list sum.
\end{itemize}

The complete multiplication proof becomes
\begin{Lean}
  simp only [interpret_eq_sum, multiply, List.flatMap_def,
    List.map_flatten, List.sum_flatten, List.map_map,
    Function.comp_def, exponent_add,
    ← MvPolynomial.monomial_mul,
    List.sum_map_mul_left, List.sum_map_mul_right]
\end{Lean}

The private left-multiplication induction is no longer needed. The structural
induction proving \code{interpret_expand} remains: it is the correct connection
between expression constructors and the established list algebra.

A small companion derives \code{exponent_inj} from the injectivity of
\code{Finsupp.equivFunOnFinite.symm}. It replaces a local extensionality proof,
and is counted within this module change rather than as a separate proposal.

\subsection{Why the executable behavior stays exact}

Replacing the list by a finite set would discard duplicate occurrences.
This refactor does not do that. Empty lists, repeated exponent vectors,
negative coefficients, and coefficients that later cancel are still processed
by the original list operations.

Even the existing recursive \code{interpret} definition is retained, although
it is used only for mathematical interpretation. Equation (5) is a theorem
about it, so its definitional behavior is preserved.

\subsection{Measurements and checking scope}

The production candidate is 110 physical lines instead of 120, and 87
nonblank lines instead of 97. These counts include the new private theorem
and the explicit list-algebra import. Five list inductions become one; the
expression induction remains unchanged.

Both complete modules compile with Lean 4.32.0, using a freshly rebuilt
current \code{Representation} module. Both public
\code{coefficient_expand} endpoints report
\[
 \{\text{\texttt{propext}},\ \text{\texttt{Classical.choice}},\ \text{\texttt{Quot.sound}}\}.
\]
There is no downstream-build or speedup claim. The compiler times in the logs
are execution receipts from a shared environment, not benchmarks.

The exact patch is \code{patches/syntax-coefficients.patch}. Complete before
and after sources are also included.

\section{P3: transport quantification by restriction surjectivity}

\subsection{The actual environment-extension problem}

\src{ObjectPolynomial}{\code{sat_noZeroForm_iff}} relates validity of a sealed
no-zero formula over environments $\N\to V$ to a predicate over finite tuples
$\operatorname{Fin}(n)\to V$. Its original proof constructs an environment by
hand: use the tuple below $n$, and a default object elsewhere. It then proves
restriction recovers the tuple and transports guards and equation
satisfaction separately.

The content is simply surjectivity of restriction. For an injection
$j:I\hookrightarrow J$ and a default $d\in V$, the map
\[
 \rho:(J\to V)\longrightarrow(I\to V),\qquad
 \rho(e)=e\circ j
\]
is surjective. Thus every predicate $P$ on assignments satisfies
\[
       (\forall e,\ P(e\circ j))
       \quad\Longleftrightarrow\quad
       (\forall a,\ P(a)).
 \tag{6}
\]
Include the naturalness guard and unsolvability conclusion in $P$.

\subsection{Use the existing extension principle}

The checked generic mechanism uses Mathlib's
\code{Function.Injective.surjective_comp_right'} and
\code{Function.Surjective.forall}:

\begin{Lean}
theorem forall_assignments
    (coordinates : ι → I)
    (hinj : Function.Injective coordinates)
    (default : V) (P : (ι → V) → Prop) :
    (∀ env : I → V, P (env ∘ coordinates)) ↔
      ∀ assignment, P assignment :=
  (hinj.surjective_comp_right' (fun _ => default)).forall.symm
\end{Lean}

For the actual theorem, the injection is \code{Fin.val} and the default is
\code{vempty}~\code{H}. No new production helper is necessary: invoke the library
fact after the existing satisfaction equivalences.

\begin{Lean}
  refine (seal_valid (mem := mem) _).trans ?_
  change (∀ ee : Nat → V,
    Sat mem ee (inputsNaturalForm (Fin.val : Fin n → Nat)) →
      ¬ Sat mem ee (equationForm e Fin.val)) ↔ NoZeroRel H e
  simp only [inputsNaturalForm_spec H, equationForm_spec H]
  exact
    (Fin.val_injective.surjective_comp_right'
      (fun _ => vempty H)).forall.symm
\end{Lean}

The displayed source is wrapped for the page; the supplied patch contains
the exact eight-line declaration.

\subsection{Boundary cases and distinctness}

Injectivity matters. If two tuple coordinates map to the same environment
coordinate, arbitrary distinct tuple values cannot be recovered. The
universal assignment conclusion would then be stronger.

The explicit default handles the empty coordinate type without silently
assuming a nonempty target type in the abstract theorem. In the model
application, the empty set already provides this value. No standardness,
arithmetic, well-foundedness, or induction property of $V$ is used.

Earlier sentence-transport proposals consume the no-zero equivalence.
P3 instead simplifies the environment extension inside that equivalence,
so the benefit remains after those earlier wrappers have been applied.

\subsection{Evidence and outstanding integration}

The generic assignment theorem, its \code{Fin.val} specialization, and the
guarded-unsolvability specialization compile. Their only audited axiom
dependencies are the usual \code{propext}, \code{Classical.choice}, and
\code{Quot.sound}.

The production declaration changes from 20 to 8 lines, a proposed
twelve-line reduction with no added helper. However, the complete
\code{ObjectPolynomial} module was not compiled: its deep object-model
dependency closure was unavailable in the local cache. The production patch
is marked \emph{integration pending}. Success of the abstract mechanism is
not presented as a complete object-module check.

\section{P4: use the order obstruction on a residual exponent}

\subsection{A smaller proof of the existing theorem}

\src{NativeBinaryScale}{\code{width_dvd_exponent}} proves
\[
      2^d-1\mid 2^u-1\quad\Longrightarrow\quad d\mid u
\]
with its retained public assumption $0<d$. Put $r=u\bmod d$. The existing
modular-power identity reduces the divisibility hypothesis to
\[
      2^r-1=0.
\]
The old proof turns this into $2^r=1$, rewrites $1$ as $2^0$, and invokes
injectivity of exponentiation.

The replacement uses the order obstruction:
\[
      2^r-1=0\Longrightarrow 2^r\le1.
\]
If $r\ne0$, \code{Nat.pow_le_one_iff} would force $2\le1$.
So $r=0$, and hence $d\mid u$.

\begin{Lean}
  have hm := Nat.mod_eq_zero_of_dvd hdiv
  rw [Nat.pow_sub_one_mod_pow_sub_one] at hm
  apply Nat.dvd_of_mod_eq_zero
  by_contra hn
  have hsmall :=
    (Nat.pow_le_one_iff hn).mp (Nat.sub_eq_zero_iff_le.mp hm)
  omega
\end{Lean}

The public signature is unchanged, including \code{hd}. The supplied source
formats the body in six lines; the page breaks one long line.
The declaration shrinks from twelve to eight lines. Both exact declarations
compile with identical axiom sets. Their unused-\code{hd} warning is also
identical.

\subsection{Optional arithmetic API: recover a block power directly}

The independently checked \code{PrimeRadix.lean} proves
\[
\begin{gathered}
 p\text{ prime},\quad 1<q,\quad q\mid p^R,\quad p^d-1\mid q-1\\
 \Longrightarrow\quad
 \exists N>0,\quad dN\le R,\quad q=(p^d)^N .
\end{gathered}
\tag{7}
\]
Classify a divisor of a prime power as $q=p^u$ with $u\le R$.
The generalized exponent-divisibility fact yields $u=dN$, and $q>1$
forces $N>0$.

An affine adapter accepts $q=1+J(p^d-1)$, matching \code{sourceQ}.
A source client can then supply its divisibility certificate and recover
the cell-count exponent and budget together, rather than exposing an
intermediate binary exponent.

The generic core needs no separate $d>0$ hypothesis: for $d=0$ the other
premises force $q=1$, contradicting $q>1$. The public source wrapper can
still keep its original parameter.

This optional API is not counted as additional removed lines. Its shared
proof cost may exceed its benefit for a single current client. Adopt it
if the arithmetic classification becomes a reused boundary.

\subsection{Necessary guards}

The $q>1$ premise is required for a positive exponent. Primality is also
substantive. For
\[
       p=6,\quad d=2,\quad q=4096,\quad R=12,
\]
one has $4096\mid6^{12}$ and $35\mid4095$, but $4096$ is not a positive power
of $36$. Thus (7) cannot be generalized to arbitrary bases larger than one.
Exponent divisibility alone needs only $p>1$; primality is needed to
classify divisors of $p^R$.

The exact declaration checks and the optional general prototype are
included. The modified complete \code{NativeBinaryScale} module and its
downstream clients were not rebuilt.

\clearpage
\section{Validation, costs, and adoption gates}

\begin{longtable}{L{0.20\linewidth}L{0.39\linewidth}L{0.30\linewidth}}
\toprule Artifact & What was checked & What remains\\ \midrule
P1 core & Actual syntax, split coefficient identity, injective coefficient
map, semiring evaluation, and cancellation. Current lower dependencies rebuilt. &
Production wrapper and full internal-model clients.\\
P1 clients & Literal conditional row, exact input slots, and Boolean
noncancellation counterexample. &
Complete \code{ObjectConditional} integration.\\
P2 & Complete original and patched sparse-coefficient modules;
current representation dependency rebuilt; endpoint axioms matched. &
Downstream clients and project checks.\\
P3 & Generic restriction theorem, finite specialization, guarded predicate. &
Complete object-polynomial elaboration.\\
P4 local & Exact old and new declarations; signature and axiom preservation. &
Enclosing module and downstream checks.\\
P4 optional & Prime-power and block-radix arithmetic statements. &
Source adapter and helper-cost comparison.\\
\bottomrule
\end{longtable}

The checked compiler is Lean 4.32.0,
commit \code{8c9756b28d64dab099da31a4c09229a9e6a2ef35}.
The repository lock records Mathlib
\code{81a5d257c8e410db227a6665ed08f64fea08e997}.
Current representation and natural-expression dependencies were rebuilt for
P1 and P2.

The supplied replay script was run successfully: four current production
dependencies, six prototype files, and two complete sparse-coefficient
modules give twelve successful compiler invocations. All three production
patches also apply to their exact recorded baseline sources; patch application
is not a replacement for the outstanding module checks.

Successful compiler receipts are included. The checked candidates contain
no \code{sorry} declarations or new axioms. The usual Lean/Mathlib axioms
appear in the semantic and coefficient endpoints; the Boolean evaluation
counterexamples themselves have no axiom dependencies.

No timing is treated as a speedup, and no finite sample replaces a universal
proof. Source-size figures are physical and nonblank lines, not a measure of
semantic complexity.

\subsection{Before adoption}

Use a disposable complete repository checkout with the pinned toolchain.
Check each patch against its recorded baseline bytes, rebuild the affected
module and direct clients, and inspect relevant endpoint axioms.

For P1, pilot a small group of algebraic row proofs. Keep integer-polynomial
equality an explicit certificate; avoid unrestricted polynomial expansion
as a global simplification rule. Preserve the literal expression, witness
data, degree claims, and model assumptions.

Retain duplicate-preserving list operations (P2), the coordinate injection
and default object (P3), and the public signature and optional classifier
assumptions (P4).

\section{Archive guide and reproduction}

The packet is a focused research deliverable, not a repository mirror:

\begin{itemize}
\item \code{report/}: this PDF and its complete editable LaTeX source.
\item \code{lean/}: checked semantic kernels, the exact conditional client,
Boolean negative control, and exponent before/after checks.
\item \code{patches/}: production edits for sparse interpretation, environment
restriction, and exponent divisibility.
\item \code{baseline/}, \code{candidates/}: small source snapshots for review.
\item \code{evidence/}: compiler receipts, line counts, source identities,
closing GitHub comparison, and exclusion evidence.
\item \code{scripts/}: static archive verification and compiler replay for a
complete configured checkout.
\end{itemize}

Replay does not install a toolchain or silently modify the original
repository. The standalone source checks and optional production-patch
integration are separate steps. The object-model replacement remains
explicitly uncompiled in this delivery.

\section*{Primary sources}
\addcontentsline{toc}{section}{Primary sources}

\begin{enumerate}
\item \href{\repo/tree/\pin/Computability/DiophantineConsistency}
{Maintained DiophantineConsistency subtree}, closing pin \code{7df43cdfacb4}.
\item \src{Representation}{integer syntax and polynomial interpretation};
\src{NaturalExpression}{natural positive/negative syntax};
\src{NaturalExpressionSemiring}{semiring evaluation}.
\item \src{GatesSemiring}{signed zero semantics};
\src{ObjectConditional}{literal selection-row semantics}.
\item \src{SyntaxCoefficients}{executable expansion and coefficient correctness}.
\item \src{ObjectPolynomial}{sealed no-zero semantics on finite internal tuples}.
\item \src{NativeBinaryScale}{binary and block-radix recovery}.
\item \href{\repo/tree/\pin/Computability/DiophantineConsistency/research/lean-proof-simplification}
{Placed proposal corpus a--w}, including substantive supplementary notes.
\item \href{\repo/tree/\pin/docs/incoming}{Five incoming ZIPs}, including
the final additional interfaces report.
\item \href{https://github.com/leanprover-community/mathlib4/tree/81a5d257c8e410db227a6665ed08f64fea08e997/Mathlib}
{Mathlib at the repository lock}: function restriction, list sums,
polynomial maps, and natural-power arithmetic.
\end{enumerate}

\end{document}
